\documentclass{article}
\usepackage{amsmath}
\usepackage{amssymb}
\usepackage{fullpage}
\usepackage{enumerate}
\usepackage{xcolor}
\usepackage{bm}
\pagestyle{empty}
\usepackage{graphicx}
\graphicspath{{C:\Users\steve\Pictures}}
\begin{document}

\noindent{{\Large \bf Fall 25, 4100 Problem Set 9}

\large

\bigskip
\noindent{Please complete all problems. For any proof questions, please structure your proof similar to those from lecture.

\begin{enumerate}

\item Sketch the conic section $x_1^2+8x_1x_2+7x_2^2=1$ by finding the spectral decomposition of the matrix $A=\left[\begin{array}{cc}
1 & 4\\
4 & 7
\end{array}\right]$


\vfill

\item Let $A$ be any real matrix and prove that
\begin{enumerate}[a.]
\item If $\bm{v}$ is an eigenvector of $AA^T$, then $\bm{v}\in{\rm range}(A)\cup{\rm null}(A^T)$ (look at the cases where the corresponding e-value is 0 or nonzero)
\item If $\lambda$ is a nonzero eigenvalue of $AA^T$, then it is also an eigenvalue of $A^TA$

\end{enumerate}

\vfill

\item Let $$A=\left[
\begin{array}{cc}
1 & 1\\
-1 & 1\\
1 & 2
\end{array}\right]$$

\begin{enumerate}[a.]
\item Find the singular value decomposition of $A$
\item Find the psuedo-inverse, $A^+$ and use it to find the vector in ${\rm range}~A$ that is closest to $\left[
\begin{array}{c}
1\\
2\\
3
\end{array}\right]$

\end{enumerate}

\vfill

\item Find and sketch the image of the unit circle under the linear transformation $T:\mathbb{R}^2\rightarrow\mathbb{R}^2$ where
\begin{eqnarray*}
T\bm{e}_1&=&2\bm{e}_1+\bm{e}_2
\\ T\bm{e}_2&=&-3\bm{e}_1+6\bm{e}_2
\end{eqnarray*}
using the method from lecture 11-20

\item Given $n\times n$ matrix $A$, define the trace of $A$ as
$${\rm tr(A)}=\sum_{i=1}^n A_{i,i}=\langle {\bm e}_1,A{\bm e}_1\rangle+\cdots +\langle {\bm e}_n,A{\bm e}_n\rangle$$

Now letting $A\in\mathbb{C}^{n\times n}$ have singular values of $\sigma_1,\cdots,\sigma_r$, show that  

$${\rm tr}(A^*A)=\sigma_1^2+\cdots+\sigma_r^2$$




\end{enumerate}

\end{document}

